\section{Preliminaries}
% In this section, before introducing the Discounted Beta–Bernoulli reward estimation, we first review RLVR and GRPO, which constitute the core prior work underpinning our approach.
We briefly introduce RLVR and its commonly used baseline, GRPO, to establish the foundation for our method.

\subsection{Reinforcement Learning with Verifiable Rewards}
Given a prompt $q \sim \mathcal{D}$ from the training dataset and a policy $\pi_{\theta}$, a generated response is $o \sim \pi_{\theta}(\cdot \mid q)$.
In RLVR, the reward function $r(\cdot)$ is typically defined as a binary signal indicating whether the response contains a correct answer.

% \vspace{-2mm}
\begin{equation}
r(o, q) =
\begin{cases}
1, & \text{if $o$ contains the answer of $q$} \\
0, & \text{otherwise}.
\end{cases}
\end{equation}
% \vspace{-2mm}

\subsection{Group Relative Policy Optimization (GRPO)}
GRPO \cite{shao2024deepseekmath} is a foundational baseline for RLVR that eliminates the need for an explicit value model and Generalized Advantage Estimation \cite{ppo}. Instead, it relies on group-relative normalization to estimate advantages. 

For a given prompt $q$, GRPO samples $N$ independent responses $\{o_i\}_{i=1}^{N}$ from the old policy $\pi_{\theta_{\mathrm{old}}}$ and computes advantages by normalizing the corresponding rewards $\{r(o_i, q)\}_{i=1}^{N}$ within the group:

% \vspace{-2mm}
\begin{equation}
\hat{A}_i
=
\frac{r(o_i, q) - \mu}{\sigma},
\end{equation}
% \vspace{-2mm}

where $\mu = \frac{\sum_{i=1}^{N} r(o_i, q)}{N}$ and $\sigma^2 = \frac{\sum_{i=1}^{N} \left(r(o_i, q) - \mu\right)^2}{N-1}$.
The resulting advantage $\hat{A}_i$ is broadcast to all tokens in the response, i.e., $\hat{A}_{i,t} = \hat{A}_i$ for all token positions $t$.

The policy is then updated by maximizing a clipped surrogate objective based on estimated advantages:
\begin{equation}
\label{eq:grpo}
\begin{aligned}
&\mathcal{J}_{\mathrm{GRPO}}(\theta)
=
\mathbb{E}_{q \sim \mathcal{D},\, \{o_i\}_{i=1}^{N} \sim \pi_{\theta_{\mathrm{old}}}(\cdot \mid q)}
\Bigg[
\frac{1}{N}
\sum_{i=1}^{N}
\frac{1}{|o_i|}
\sum_{t=1}^{|o_i|}\\
&\min \Big(
w_{i,t}(\theta)\, \hat{A}_{i,t},
\mathrm{clip}\!\big(w_{i,t}(\theta), 1-\epsilon, 1+\epsilon\big)\, \hat{A}_{i,t}
\Big)
\Bigg],
\end{aligned}
\end{equation}
where $\epsilon$ is the clipping hyperparameter and 
$w_{i,t}(\theta)=\frac{\pi_\theta(o_{i,t} \mid q, o_{i,<t})}{\pi_{\theta_{\mathrm{old}}}(o_{i,t} \mid q, o_{i,<t})}$ denotes the per-token importance weight. Following DAPO \cite{yu2025dapoopensourcellmreinforcement}, we omit the KL divergence term between the online and reference policies, as it may restrict exploration.
